\section{Slice Properties}
\label{sec:slice}

Four constructions for extracting a sublist by index range --- all
equivalent.

\begin{itemize}
  \item \href{https://github.com/thiagomata/prime-numbers/blob/list-article-v1.0.0/src/main/scala/v1/chapter3/list/ListUtils.scala}{Tail-recursive slice}:
    builds from the end by prepending
  \item \href{https://github.com/thiagomata/prime-numbers/blob/list-article-v1.0.0/src/main/scala/v1/chapter3/list/properties/SliceEquivalenceLemmas.scala}{Head-recursive slice}:
    builds from the front
  \item \href{https://github.com/thiagomata/prime-numbers/blob/list-article-v1.0.0/src/main/scala/v1/chapter3/list/properties/SliceEquivalenceLemmas.scala}{Index-range slice}:
    accumulates by position within a range
  \item \href{https://github.com/thiagomata/prime-numbers/blob/list-article-v1.0.0/src/main/scala/v1/chapter3/list/properties/ListUtilsProperties.scala}{Slice append consistency}:
    appending a singleton preserves the slice structure
\end{itemize}

\subsection{Tail-Recursive Slice}

The tail-recursive slice builds the sublist from the end, prepending
elements as it recurses backward.

\begin{equation*}
\forall\, L \in \mathbb{L}, \forall\, i, j \in \mathbb{N},\ i \leq j < |L|
\end{equation*}

\begin{equation*}
\operatorname{slice}(L, i, j) :=
\begin{cases}
L_j :: L_e & \text{if } i = j \\
\operatorname{slice}(L, i, j - 1) \mathbin{\texttt{++}} (L_j :: L_e) & \text{if } i < j
\end{cases}
\end{equation*}

\textbf{Goal}:

\begin{equation*}
\forall\, L \in \mathbb{L}, \forall\, i, j \in \mathbb{N},\ i \leq j < |L| \implies \operatorname{slice}(L, i, j) = L[i \dots j]
\end{equation*}

\textbf{Proof by induction on $j$, with fixed $i$}

\textbf{Base case}: $j = i$

\begin{equation*}
\operatorname{slice}(L, i, i) = L_i :: L_e = L[i \dots i]
\end{equation*}

\textbf{Inductive step}: Assume

\begin{equation*}
\operatorname{slice}(L, i, j - 1) = [ L_k \mid i \leq k \leq j - 1 ]
\end{equation*}

Show:

\begin{equation*}
\begin{aligned}
\operatorname{slice}(L, i, j) &= \operatorname{slice}(L, i, j - 1) \mathbin{\texttt{++}} (L_j :: L_e) && \text{[by definition of slice]} \\
&= L[i \dots (j - 1)] \mathbin{\texttt{++}} (L_j :: L_e) && \text{[by Inductive Hypothesis]} \\
&= [ L_k \mid i \leq k \leq j - 1 ] \mathbin{\texttt{++}} (L_j :: L_e) && \text{[by Specification]} \\
&= [ L_k \mid i \leq k \leq j ] && \text{[by definition of Concatenation]} \\
&= L[i \dots j] && \text{[Q.E.D]} \\
\end{aligned}
\end{equation*}

\begin{equation*}
\therefore
\end{equation*}

\begin{equation*}
\forall\, 0 \leq i \leq j < |L|,\ \operatorname{slice}(L, i, j) = L[i \dots j]
\quad \blacksquare
\end{equation*}

This property is verified in
\href{https://github.com/thiagomata/prime-numbers/blob/list-article-v1.0.0/src/main/scala/v1/chapter3/list/ListUtils.scala}{\texttt{ListUtils::slice}}.
The implementation and verification code are in Appendix~A.3.

\subsection{Head-Recursive Slice}

The head-recursive slice builds the sublist from the front, cons-ing
elements as it recurses forward.

\begin{equation*}
\forall\, L \in \mathbb{L}, \forall\, i, j \in \mathbb{N},\ i \leq j < |L|
\end{equation*}

\begin{equation*}
\operatorname{headRecursiveSlice}(L, i, j) :=
\begin{cases}
L_i :: L_e & \text{if } i = j \\
L_i :: \operatorname{headRecursiveSlice}(L, i + 1, j) & \text{if } i < j
\end{cases}
\end{equation*}

\textbf{Goal}:

\begin{equation*}
\forall\, L \in \mathbb{L}, \forall\, i, j \in \mathbb{N},\ i \leq j < |L| \implies \operatorname{headRecursiveSlice}(L, i, j) = L[i \dots j]
\end{equation*}

\textbf{Proof by induction on $j - i$}

\textbf{Base case}: $i = j$

\begin{equation*}
\operatorname{headRecursiveSlice}(L, i, i) = L_i :: L_e = L[i \dots i]
\end{equation*}

\textbf{Inductive step}: Assume

\begin{equation*}
\operatorname{headRecursiveSlice}(L, i + 1, j) = [ L_k \mid i + 1 \leq k \leq j ]
\end{equation*}

Show:

\begin{equation*}
\begin{aligned}
\operatorname{headRecursiveSlice}(L, i, j) &= L_i :: \operatorname{headRecursiveSlice}(L, i + 1, j) && \text{[by definition]} \\
&= L_i :: L[i + 1 \dots j] && \text{[by Inductive Hypothesis]} \\
&= [ L_k \mid i \leq k \leq j ] = L[i \dots j] && \text{[by specification]} \\
\end{aligned}
\end{equation*}

\begin{equation*}
\therefore
\end{equation*}

\begin{equation*}
\forall\, 0 \leq i \leq j < |L|,\ \operatorname{headRecursiveSlice}(L, i, j) = L[i \dots j]
\quad \blacksquare
\end{equation*}

{\raggedright
This property is verified in
\href{https://github.com/thiagomata/prime-numbers/blob/list-article-v1.0.0/src/main/scala/v1/chapter3/list/properties/SliceEquivalenceLemmas.scala}{\texttt{SliceEquivalenceLemmas::\allowbreak headRecursiveSlice}}.
The full Scala verification code is in Appendix~A.4.
\par}

\subsection{Index-Range Slice}

The index-range slice builds the sublist by direct index access,
recursing forward through the index range.

\begin{equation*}
\forall\, L \in \mathbb{L}, \forall\, i, j \in \mathbb{N},\ i \leq j < |L|
\end{equation*}

\begin{equation*}
\operatorname{indexRangeValues}(L, i, j) :=
\begin{cases}
L_i :: L_e & \text{if } i = j \\
L_i :: \operatorname{indexRangeValues}(L, i + 1, j) & \text{if } i < j
\end{cases}
\end{equation*}

\textbf{Goal}:

\begin{equation*}
\forall\, L \in \mathbb{L}, \forall\, i, j \in \mathbb{N},\ i \leq j < |L| \implies \operatorname{indexRangeValues}(L, i, j) = L[i \dots j]
\end{equation*}

\textbf{Proof by induction on $j - i$}

\textbf{Base case}: $i = j$

\begin{equation*}
\operatorname{indexRangeValues}(L, i, i) = L_i :: L_e = L[i \dots i]
\end{equation*}

\textbf{Inductive step}: Assume

\begin{equation*}
\operatorname{indexRangeValues}(L, i + 1, j) = [ L_k \mid i + 1 \leq k \leq j ]
\end{equation*}

Show:

\begin{equation*}
\begin{aligned}
\operatorname{indexRangeValues}(L, i, j) &= L_i :: \operatorname{indexRangeValues}(L, i + 1, j) && \text{[by definition]} \\
&= L_i :: L[i + 1 \dots j] && \text{[by Inductive Hypothesis]} \\
&= [ L_k \mid i \leq k \leq j ] = L[i \dots j] && \text{[by specification]} \\
\end{aligned}
\end{equation*}

\begin{equation*}
\therefore
\end{equation*}

\begin{equation*}
\forall\, 0 \leq i \leq j < |L|,\ \operatorname{indexRangeValues}(L, i, j) = L[i \dots j]
\quad \blacksquare
\end{equation*}

{\raggedright
This property is verified in
\href{https://github.com/thiagomata/prime-numbers/blob/list-article-v1.0.0/src/main/scala/v1/chapter3/list/properties/SliceEquivalenceLemmas.scala}{\texttt{SliceEquivalenceLemmas::\allowbreak indexRangeValues}}.
The full Scala verification code is in Appendix~A.5.
\par}

\subsection{Slice Append Consistency}

\textbf{Lemma:} A slice of a list from index $f$ to $t$ can be expressed
as the slice from $f$ to $t - 1$ concatenated with the element at index
$t$, for $f \le t < |L|$.

\begin{equation*}
\begin{aligned}
L[f \dots t] &= \operatorname{slice}(L, f, t) \\
             &= \operatorname{slice}(L, f, t - 1) \mathbin{\texttt{++}} (L_t :: L_e) \\
             &= L[f \dots t - 1] \mathbin{\texttt{++}} (L_t :: L_e) \quad \blacksquare
\end{aligned}
\end{equation*}

{\raggedright
This property is verified in
\href{https://github.com/thiagomata/prime-numbers/blob/list-article-v1.0.0/src/main/scala/v1/chapter3/list/properties/ListUtilsProperties.scala}{\texttt{ListUtilsProperties::\allowbreak assertAppendToSlice}}.
The full Scala verification code is in Appendix~A.6.
\par}
